
% ════════════════════════════════════════════
% 具体分析（5.X.1.具体分析.tex）写作规则 —— 模板内嵌，AI 先读本文件再写
% ════════════════════════════════════════════
% 【定位】本块为每个问题的「具体分析」（三级子块），是「问题X的模型的建立和求解」里最先出现的子块。
% 【结构】三段式纯文字 + 流程图；不要分点表达。
% 【开头段】同时满足三点：① 说明针对问题几；② 点出问题类型（评价/分类/预测/优化/机理分析）；③ 点出本问大致建模方法。
% 【中间段】按「首先、接着、然后、最后」写具体步骤，每个步骤不换段。
% 【结尾段】引出图：「具体分析如图\ref{fig:...}所示」，不需要额外说明。
% 【流程图】同总体分析风格（模板式/自定义均可，配色清晰）；每框 ≤7 字，figure 用 [ht]，禁止 [H]/[!h] 强制定位。
% 【禁止】本块不讲检验、不讲结果（还没求解）。
% ════════════════════════════════════════════
\subsubsection{具体分析}

% 三段：① 问题类型+建模方法 → ② 首先/接着/然后/最后 具体步骤 → ③ 引出流程图
针对问题X，本问属于<评价/分类/预测/优化/机理分析>类问题，采用...方法进行求解。

首先对...进行...，接着利用...方法...，然后构建...模型...，最后通过...得到...。

具体分析流程如图\ref{fig:analysisX_flow}所示。

% 【图】在此插入「问题X分析流程图」（\begin{figure}[ht] + tikz 绘制，每框 ≤7 字，figure 用 [ht] 禁 [H]）。
% 模板不预置绘图方块，避免空白编译时因未定义样式/缺图而失败。图引用 \ref{fig:analysisX_flow}，label 为 fig:analysisX_flow。
